% BHU PYQ -- Professional Exam Template
\documentclass[11pt,a4paper]{article}

\usepackage[a4paper,top=2.5cm,bottom=2.5cm,left=2.8cm,right=2.8cm,headheight=14pt]{geometry}
\usepackage{amsmath,amssymb}
\usepackage[T1]{fontenc}
\usepackage[utf8]{inputenc}
\usepackage{lmodern}
\usepackage{microtype}
\usepackage{fancyhdr}
\usepackage{enumitem}
\usepackage{hyperref}
\usepackage{lastpage}
\usepackage{parskip}

\hypersetup{colorlinks=true,urlcolor=black,linkcolor=black,
  pdftitle={B.Sc. (Hons.) Botany Sem-IV BOB-403A -- BHU PYQ}}

\pagestyle{fancy}
\fancyhf{}
\renewcommand{\headrulewidth}{0.4pt}
\renewcommand{\footrulewidth}{0.4pt}
\fancyhead[L]{\small   \textit{Banaras Hindu University}}
\fancyhead[R]{\small   \textit{Botany Paper: BOB-403A}}
\fancyfoot[C]{\small Page  \thepage\ of \pageref{LastPage}}


\newlist{parts}{enumerate}{2}
\setlist[parts,1]{label=(\alph*),leftmargin=*,itemsep=4pt,topsep=3pt}
\setlist[parts,2]{label=(\roman*),leftmargin=*,itemsep=2pt,topsep=2pt}

\newcommand{\pts}[1]{\hfill   \textit{[#1]}}


\begin{document}


\begin{center}
  {\large\bfseries BANARAS HINDU UNIVERSITY}\\[2pt]
  {\normalsize B.Sc. (Hons.) Semester IV Examination, 2023-24}\\[6pt]
  \rule{0.8\linewidth}{0.5pt}\\[6pt]
  {\large\bfseries Botany}\\[4pt]
  {\large\bfseries Paper: BOB-403A --- Ancillary Botany-II}\\[4pt]
  {\normalsize Time: Three Hours \hfill Maximum Marks: 70}
\end{center}

\rule{\linewidth}{0.4pt}

\smallskip



\noindent   \textit{Instructions: Answer five questions including Question 1 which is compulsory. The figures in the right-hand margin indicate marks.}

\bigskip

\medskip


\noindent   \textbf{Question 1.} Fill in the blanks or answer the following questions (in not more than 50 words each):\pts{$10  \times 1 = 10$}

\begin{parts}
  \item What is substrate-level phosphorylation?
  \item What is proto-cooperation?
  \item Define apical dominance.
  \item What are hydrophytes?
  \item Gibberellins were first recognised from the fungus \makebox[3cm]{\hrulefill}.
  \item What are the main constituents of   \textit{Tinospora cordifolia}?
  \item Define Endemic disease.
  \item What is the role of microconsumers in the ecosystem?
  \item The disease Citrus canker is caused by \makebox[3cm]{\hrulefill}.
  \item Why are nutrient cycles also referred to as biogeochemical cycles?
\end{parts}

\medskip


\noindent   \textbf{Question 2.} What are plant growth hormones? Describe the various roles of ethylene in detail in plants.\pts{15}

\medskip


\noindent   \textbf{Question 3.} What are biotic interactions in the ecosystem? Describe different kinds of possible negative interactions among species.\pts{15}

\medskip


\noindent   \textbf{Question 4.} Discuss the structural components and function of the ecosystem in detail. What is the importance of the food chain in the ecosystem?\pts{15}

\medskip


\noindent   \textbf{Question 5.} Discuss how nutrient cycling takes place in the ecosystem. Explain the steps of the carbon cycle in the atmosphere.\pts{15}

\medskip


\noindent   \textbf{Question 6.} Write short notes on any two of the following:\pts{$2  \times 7.5 = 15$}

\begin{parts}
  \item Food pyramid
  \item Disease symptoms and mode of spread and transmission of Papaya mosaic virus
  \item Sulphur cycle
\end{parts}

\medskip


\noindent   \textbf{Question 7.} Write short notes on any two of the following:\pts{$2  \times 7.5 = 15$}

\begin{parts}
  \item Medicinal uses and chemical constituents of   \textit{Withania somnifera}
  \item Mutualism with suitable examples
  \item Plant responses to changes in temperature
\end{parts}

\medskip


\noindent   \textbf{Question 8.} Answer/write short notes on any three of the following:\pts{$3  \times 5 = 15$}

\begin{parts}
  \item What is photoperiodism? Classify plants based on light period requirements.
  \item Phloem loading
  \item Structure of phytochrome
  \item Electron transport system in mitochondria
\end{parts}

\medskip


\noindent   \textbf{Question 9.} Define plant tissue culture. Describe in detail the various techniques of plant tissue culture along with their applications.\pts{15}

\vfill
\rule{\linewidth}{0.4pt}

\begin{center}\small
\href{https://bhu-exam.vercel.app/}{Click Here}\\[4pt]\href{https://me-aryan.vercel.app/}{Click Here}\\[4pt]\href{https://quashan-me.vercel.app/}{Click Here}
\end{center}

\end{document}
